\paragraph{A fixed helper alphabet.}
Unbounded owner alternation does not require unbounded repetition of the
same helper. Fix $K\geq3$. A helper letter $e=(r_e,H_e)$ acts on the current value by
$A\mapsto A\star_{r_e}H_e(w)$, where $H_e$ omits the distinguished input
$X$. For a word $u$ in these letters, write $S_u(X,w)$ for its spine.
The helpers may be arbitrary nonempty-valued operations in the value
statements below; the marked-type statements assume they are core.

\begin{theorem}[Last-occurrence bounds and triangular blocks]
\label{core:thm:spine-block-normal-form}
Suppose a spine word uses $m$ distinct helper letters. Order them
$f_1,\ldots,f_m$ by their last occurrence, and write
\[
 P=f_1f_2\cdots f_m,\qquad
 T=(f_m\cdots f_1)(f_m\cdots f_2)\cdots(f_m).
\]
For $m=0$ take both words empty.
Then
\[
 S_P\subseteq S_u\subseteq S_T.
\]
All three have the same singleton-output tests. For core helpers they
also have the same exact marked full type, which depends only on the
used alphabet. Moreover, $u$ has an equivalent subword of length at most
$m(m+1)/2$.
\end{theorem}
\begin{proof}
Fix a passive valuation. Each letter is either a constant reset, when
$H_e=\{r_e\}$, or a monotone extensive map. Deleting a nonfinal copy of
a letter therefore decreases the final value: an extensive copy can only
enlarge its input, while a resetting copy is made irrelevant by its later
copy. At any singleton-output test, a letter acts on its Boolean bit by
identity or by a constant. Its final copy makes every earlier copy
irrelevant. Deleting all nonfinal copies gives $P$ and preserves these
tests. The marked factors are conjunctions of the helper filters, so
repeating a letter does not change them.

Cut $u$ into blocks, block $i$ ending at the last $f_i$.
Its letters belong to $\{f_i,\ldots,f_m\}$ and it ends in $f_i$.
We show that the value depends only on these block contents and final
letters. If resets occur, let $h$ be the largest last-occurrence index of
a resetting letter. The reset at the end of block $h$ erases the prefix;
put $p=r_{f_h}$. All later letters are nonresetting. If there is no reset
and the input is an atom, put $h=0$ and let $p$ be that atom's owner.
A letter fixes $\{p\}$ exactly when its owner is $p$ or its helper is
$\{p\}$.

Let $\beta$ be the first block after $h$ containing a letter not fixing
$\{p\}$. If none exists, the final value is $\{p\}$. Otherwise the
value becomes ambiguous in that block. The final $f_\beta$ lies at or after
this change, and every letter with later last index still has a copy ahead.
Thus the final value is exactly
\[
 \{p\}\cup\bigcup_{j\geq\beta}H_{f_j}(w).
\]
This is independent of repetitions and internal order within each block,
provided $f_i$ stays last. With no reset and an ambiguous input, every
gate is union and the same independence is immediate.

Keep one copy of each letter in each block, retaining $f_i$ at its end.
This gives an equivalent subword with at most $m-i+1$ letters in block $i$,
hence the triangular bound. To obtain the upper bound, insert every missing
later letter into each block. Each insertion is before that letter's final
copy and can only increase the context, by the deletion argument.
Block-order independence identifies the result with $T$.
\end{proof}

\begin{corollary}[Safety reduces to a permutation]
\label{core:cor:spine-permutation-safety}
For $K\geq3$ and a specified finite helper alphabet, a spine using every letter and
satisfying $S_u\subseteq C\cup X$ exists if and only if a permutation
of that alphabet satisfies the same bound. Singleton-profile constraints
and the exact marked full type can be preserved as well.

For a permutation $f_1,\ldots,f_m$, the following is an exact safety test
at each passive valuation. Put $A_i=H_{f_i}(w)$, and let $h$ be the largest
index with $A_h=\{r_{f_h}\}$, or zero if there is no such index.
\begin{enumerate}
\item If $h=0$, require $A_i\subseteq C(w)$ for every $i$.
\item If $h>0$, put $p=r_{f_h}$ and require $p\in C(w)$. For each
  $k\notin C(w)$ let $\ell_k=\max\{i:k\in A_i\}$, with maximum zero
  for the empty set. Require
  \[
   r_{f_i}=p\ \text{or}\ A_i=\{p\}
   \qquad(h<i\leq\ell_k).
  \]
\end{enumerate}
\end{corollary}
\begin{proof}
The least permutation in \cref{core:thm:spine-block-normal-form} proves
the existence equivalence. With no reset, an ambiguous input
$X=I\setminus\{k\}$ detects each forbidden helper winner $k$, so all
helpers must be bounded by $C$; this also suffices because every gate is
contained in its operands' union. After a last reset, the output is
independent of $X$, so its bound for every $X$ is containment in $C$.
It always contains $p$. A forbidden $k$ appears exactly when the first change from $\{p\}$
occurs no later than its last helper contribution, namely $\ell_k$.
The stated fixing conditions exclude exactly that event.
\end{proof}

Exact contexts can still need quadratically many copies of their letters.
This concerns equivalence over a fixed elementary alphabet.

\begin{proposition}[Quadratic length is necessary for exact contexts]
\label{core:prop:spine-quadratic-alphabet}
Let $s,t\geq1$ and choose distinct owners $a,b,c$.
Use independent helper variables $x_1,\ldots,x_s$ at owner $a$ and
$y_1,\ldots,y_t$ at owner $b$. The word
\[
 (y_1\cdots y_t)x_1\,(y_1\cdots y_t)x_2\,\cdots\,
 (y_1\cdots y_t)x_s\,(y_1\cdots y_t)
\]
has length $st+s+t$, and no shorter word over the same elementary
alphabet has its unary context on all winner-valued inputs.
In particular, for $m=s+t$ letters the worst-case length is at least
$\lfloor m^2/4\rfloor+m$.
\end{proposition}
\begin{proof}
Every letter must occur in an equivalent word: make it alone reset to
its owner's atom, set every other helper and $X$ to $\{c\}$, and use
provenance to detect its owner. Write $\operatorname{last}(e)$ for a
letter's last position in such a word.

Set $x_i=\{a\}$, $y_j=\{b\}$, and every other helper to $\{c\}$.
The output retains the owner of the last reset and cannot acquire the
other reset owner afterwards. In the displayed word the last reset is
$y_j$, so equivalence forces
$\operatorname{last}(x_i)<\operatorname{last}(y_j)$ for every $i,j$.

For $2\leq i\leq s$, set $X=\{a\}$,
$x_{i-1}=\{a\}$, $x_i=\{b\}$, $y_j=\{c\}$, every other $x$ to
$\{c\}$, and every other $y$ to $\{a\}$.
There is just one resetting letter, $x_{i-1}$. Afterwards the only
letter that can first change the atom $a$ is $y_j$, and the only source
of $b$ is $x_i$, whose own $a$-gate fixes that atom until it changes.
The displayed word includes $b$, so every equivalent word has a copy of
$y_j$ after the last $x_{i-1}$ and before some $x_i$.
In particular the last $x_i$ positions strictly increase, and every
consecutive gap contains a copy of every $y_j$.
For the first gap use the same test without a reset: $X=\{a\}$,
$x_1=\{b\}$, $y_j=\{c\}$, other $x$ equal to $\{c\}$ and other $y$
equal to $\{a\}$. It forces $y_j$ before the last $x_1$.
Finally the reset-order test forces another $y_j$ after the last $x_s$.
These $s+1$ disjoint gaps require $t(s+1)$ occurrences of $y$-letters,
in addition to at least $s$ occurrences of $x$-letters.
The displayed word attains this lower bound. Balancing $s,t$ gives the
last assertion.
\end{proof}
